\scanpage{431}
\mspage{218}

il en est de $\overline{\mathrm m}$ \unsure{pour} l'action induite sur chaque $\pi(I')$ $(I'\subset I=I(\pi))$ [il suffit de prendre les $\pi(I')$ pour $I'$ de la forme $I\setminus\{j\}$, $j\in I$] --- \uuline{et} qu'enfin l'action de $\Gamma^{\natural}$ \deleted{relatif à \unsure{un} $i\in I$} sur $\pi'(I)$ (déduite de $\pi'$ en bouchant les trous en $I$, de sorte que $\pi'(I)$ a une seule classe de lacets; \unsure{$\pi'(I)$} se déduit aussi de l'action \inserted{effective} de $\Gamma$ sur $\pi(I)$ --- avec $0$ \inserted{classe de} lacets --- en faisant ``un trou'' correspondant) relatif à \inserted{un} point $i\in I$, ne normalise pas de ss-groupes lacets\ldots\ Mais je présume que la première condition (action \inserted{de $\Gamma^{\natural}$} ne normalisant aucun des $L_i$) implique les autres.

\margin{du moins si $g\ne 0$}

Mais il faut \inserted{dans la définition d'admissibilité} aussi tenir compte de la condition $(4^\circ)$, \deleted{\illegible{1}}\inserted{qui implique l'}existence de suffisamment de relèvements\ldots\ .

J'en arrive (péniblement!) à une

\uuline{Conjecture provisoire profinie} Pour $\Gamma$ groupe profini donné, une action extérieure sur un $(\pi,a)$ arithmétisé anabélien de type $(g,\nu)$ est dite \uline{admissible}, si

\margin{Canulé cf plus bas\ldots}

a) \deleted{\illegible{1}} L'hom $\Gamma\longrightarrow\boldsymbol{\Gamma}_a$ a une image ouverte

b) Les \unsure{actions} \uline{effectives} déduites par bouchages \unsure{de|des} trous $(i\in I'\subset I=I(\pi))$ ne normalisent pas de ss-groupes lacets.

\scanpage{432}

c) Il existe une infinité de \inserted{classes de $\pi$-conj. des} germes \deleted{\illegible{1}} scindages de l'extension de $\Gamma$ par $\pi$, qui ne normalisent aucun ss-groupe à lacets de $\pi$.

\margin{Cette \inserted{condition c)} exclut sans doute des cas comme $\Gamma=\unsure{\mathcal N_{g,\nu}|N_{g,\nu}}$ avec $(g,\nu)\ne(0,3)$, car l'extension $\unsure{\mathcal M_{g,\nu}|M_{g,\nu}}$ de $\unsure{\mathcal N_{g,\nu}|N_{g,\nu}}$ par $\unsure{\widehat{\pi}_{g,\nu}|\widehat{\Pi}_{g,\nu}}$ n'admet sans doute pas de germes de \illegible{1} ($\overset{?}{\Longleftrightarrow}$ $\unsure{\mathcal U_{g,\nu}|U_{g,\nu}}$ sur $\unsure{\mathcal M_{g,\nu}|\mathcal N_{g,\nu}}$ \unsure{n'a|n'admet} sans doute pas de \unsure{multisection étale|multisections étales}\ldots)}

Une action effective de $\Gamma$ sur $\pi$ est dite effective, \deleted{et} si l'action ext. est effective, et si elle ne normalise aucun ss-groupe à lacets.

Ceci posé:

a) Si $\Gamma$ opère effectivement de façon admissible, on a
\[
\pi^\Gamma=(1)
\]
(peut-être $\overline{\mathrm m}$ $\pi^{\unsure{\Gamma^+|\Gamma^\dagger}}=(1)$) \deleted{je présume aussi que les seuls \unsure{sous-}}

\deleted{b) Si $\Gamma$ opère extérieurement de façon admissible, \deleted{alors} en laissant stable $I'\subset I(\pi)$, alors itou sur $\pi(I')$ --- de $\overline{\mathrm m}$ pour actions effectives admissibles.}

\margin{\deleted{\illegible{1}}}

\deleted{c) Sous les conditions précédentes,}

\scanpage{433}
\mspage{219}

\deleted{l'application canonique des $\Gamma$-ensembles
\[
\mathrm{Pt}(B_{\pi,\Gamma^{\natural}})\longrightarrow
\mathrm{Pt}(B_{\pi(I'),\Gamma^{\natural}})
\]}

b) Si $i\in I$ est stable par $\Gamma$ opérant \deleted{effective}\inserted{extérieure}ment de façon admissible, alors l'action effective sur $\pi'=\pi(i)$ ne normalise aucun $L'_j$ $(j\in I\setminus\{i\})$.

c) Si $\Gamma$ opère effectivement sur $\pi$ de façon admissible, \inserted{\unsure{en} laissant fixe $i\in I$,} alors l'action correspondante effective (\uline{par} \deleted{\illegible{3}}\inserted{passage au quotient}) sur $\pi^{\prime\natural}=\unsure{\pi(i)|R(i)}$ ne normalise aucun $L'_j$ $(j\in I\setminus\{i\})$.

d) Si $\Gamma$ \deleted{opère effectivement sur $\pi$ de façon}\inserted{opère effectivement sur $\pi$ de} façon admissible, alors l'application
\[
\unsure{\beta}\qquad
\mathrm{Pt}_{\mathrm{ad}}(B_{\pi,\Gamma^{\natural}})\longrightarrow
\mathrm{Pt}_{\mathrm{ad}}(B_{\pi',\Gamma^{\natural}})
\]
définie par c) est \uline{injective}, et \inserted{le complémentaire de} son image est \deleted{finie} égale à $\{P\}$, où $P$ est défini par le ``trou'' $i$.

Mais cette dernière partie de l'assertion, allant au delà de la seule \uline{injectivité} --- et caractérisant l'image, me semble maintenant tout à fait douteuse ---

\scanpage{434}

en effet, il suffit de regarder un schéma $S$ de paramètres, de type fini sur \unsure{$Q$|$\mathbb{Q}$} (un $K(\pi,1)$ de préférence, p.ex une courbe \unsure{alg.}) et de prendre pour $\Gamma$ (non le groupe fondamental de son point générique, \inserted{i.e. d'un corps,} \unsure{mais)} le groupe fondamental de $S$ lui-$\overline{\mathrm m}$. La donnée d'une action extérieure admissible de $\Gamma$ sur un $\pi$ \unsure{correspond.} moralement à celle d'une famille de courbes $U$ de type $g,\nu$ paramétré par $S$. La condition c) est vérifiée p.ex. si $U$ provient d'une courbe alg. sur le corps de base lui-$\overline{\mathrm m}$, donc pas de pb -- et il est manifeste que la surjectivité déconne. La difficulté provient du fait que deux sections \uline{distinctes} de $U$ sur $S$ peuvent ne pas être \uline{disjointes} !

\margin{Les scindages d'extensions correspondent \unsure{aux} sections de $U$ sur $S$}

\scanpage{435}
\mspage{220}

Il faut pouvoir parler de sections ``\unsure{strictmt|strictement} distinctes'' i.e. distinctes en tt pt de $S$ -- ce qui, en traduction profinie, \unsure{revient} à comparer deux scindages de l'ext. de $\Gamma$ par $\pi$, avec \deleted{\unsure{deux}}\inserted{les} \inserted{germes de} scindages sur \unsure{$\boldsymbol{\Gamma}'_{\cdot}\subset\boldsymbol{\Gamma}$|$\boldsymbol{\Gamma}'\subset\boldsymbol{\Gamma}$} de $\Gamma$ (en tant \unsure{qu'extension} de $\boldsymbol{\Gamma}_Q$ par une partie géométrique) qui correspondent aux \unsure{pts} de $S$ -- i.e. les (germes de scindages) \uline{admissibles} justement. On a l'impression de tourner dans un cercle vicieux ou presque -- il semblerait qu'il ne faudrait pas trop vouloir avaler à la fois -- traiter d'emblée \uline{tous} les groupes profinis $\Gamma$ à la fois -- mais plutôt \unsure{se} borner \inserted{d'abord} à ceux qui, moralement, correspondent à des \inserted{\unsure{$\pi_1$|$\Pi_1$} de} schémas de type fini sur \unsure{$Q$|$\mathbb{Q}$} -- des $K(\pi,1)$ disons, ou \inserted{$\overline{\mathrm m}$} des variétés

\scanpage{436}

élémentaires à la Artin -- et dans les définitions resp. conjectures, se tirer par les lacets de souliers, en récurrant sur la dimension.

Au premier cran donc (dim $0$) on se bornerait à des actions de groupes $\Gamma$ qui soient (modulo tt au moins passage à un ss-groupe ouvert) ``arithmétiquement fidèles'' \unsure{p.e.|p.ex.} $\Gamma\to\boldsymbol{\Gamma}_a$ injectif.

\unsure{Dans|\uline{Dans}} ce cas-là, la conjecture telle qu'elle est énoncée tantôt semble raisonnable, ou du moins pas nécessairement déconnante. Le test décisif\unsure{,|,:} il est vrai, serait la possibilité d'obtenir les ``courbes'' ainsi définies comme des revêtements de $\mathbb{P}^1\setminus\{0,1,\infty\}$ \ldots

\scanpage{437}
\mspage{221}

\uline{33} Digression topologique :
Antiinvolutions des surfaces orientées \unsure{algébroïdes}

\margin{Finalement je ne regarde que les surfaces compactes -- Pourtant le cas non comp. serait aussi très intéressant à regarder -- pour essayer de retrouver par voie ``algébrique'', sur l'auto ext. d'ordre $2$ dans $\unsure{\pi_1}$, la disposition des ``points à l'infini'' de la courbe sur les composantes connexes de \deleted{\illegible{1}} $X^\sigma$ $(X=\widehat U)$}

\unsure{On appelle} surface \unsure{algébroïde} une surface $U$ de la forme $\widehat X\setminus S$, $X$ \inserted{surface} compacte orientable, $S$ fini. Alors $X$ est déterminé à homéom. unique près comme ``compactifié pur'' de $U$, \unsure{le note|on le note} $\widehat U$. \deleted{\unsure{Les}} Les homéomorphismes $U$ sur lui-$\overline{\mathrm m}$ s'identifient aux hom. de $X$ qui appliquent $S$ dans lui-$\overline{\mathrm m}$. Les orientations de $X$ sont en cor. $1$-$1$ avec celles de $U$. Les antiinvolutions de $U$ (\unsure{supposée} orientée) sont en cor. $1$-$1$ avec celles de $X$ qui conservent $S$.

On trouve alors une équivalence entre la cat. des surfaces orientées \unsure{algébroïdes} \inserted{$U$,} munies d'une antiinvolution $\sigma$, et des surfaces à bord compactes $Y$, munies d'une partie finie $T$, à $(U,\sigma)$ correspond $(\widehat U/\sigma,(\widehat U\setminus U)/\sigma)$ -- et à $(Y,T)$ correspond $\widetilde Y-\widetilde T$, où $\widetilde Y$ est le ``double orienté'' de $Y$ \inserted{(qui est une surface orientée \unsure{cpte})}, et $\widetilde T$ l'image inverse de $T$ dans $\widetilde Y$.

Nous nous intéressons maintenant au

\scanpage{438}

cas où $U$ est connexe de type $(g,\nu)$, en examinant d'abord le cas $\nu=0$, \inserted{i.e. $U=\widehat U$, $S=\varnothing$} (auquel le cas général \unsure{se ramènera}). \deleted{Cela} \inserted{La condition $\nu=0$} correspond au cas où $T=\varnothing$ -- donc les \inserted{$U=X$} \deleted{$X$} envisagés \deleted{sont} s'identifient aux doublements orientés de certaines variétés à bord compactes $Y$. Lesquelles ? Évidemment il faut que $Y$ soit \uline{connexe}, et \uuline{non orientable}.

\margin{ce n'est pas vrai que $Y$ soit nécessairement non orientable -- seulement dans le cas où $X^\sigma=\varnothing$}

Donc $Y$ est \unsure{caractérisé}, à homéomorphisme près, par son genre $\gamma$ \inserted{et le nombre $\mu$ des composantes connexes \inserted{du bord}} -- il peut se déduire du plan projectif \inserted{réel $Y_0$} en y découpant $\gamma$ rondelles disjointes, et y recollant des rubans de Möbius, \inserted{puis en \unsure{réenlevant} $\mu$ rondelles ouvertes,} -- ce qui donne
\[
\begin{aligned}
\chi(Y)&=\underset{\substack{\text{''}\\1}}{\chi(Y_0)}
-\deleted{\unsure{\chi}}\unsure{\gamma}\cdot
\underbrace{\chi_!(\text{rondelle ouverte})}_{+1}
+\unsure{\gamma}\chi_!(\text{ruban de Möbius ouvert})\\
&\hspace{11em}\unsure{+|-}\mu\underbrace{\chi_!(\text{rondelle ouverte})}_{+1}
\end{aligned}
\]
or le ruban de Möbius \inserted{ouvert} (déduit de $Y_0$ en \unsure{enlevant} une rondelle \inserted{fermée}) a \unsure{un|une} $\chi_!$ ég. à :
\[
\underset{\substack{\text{''}\\1}}{\chi(Y_0)}-
\underset{\substack{\text{''}\\1}}{\chi(\text{rondelle fermée})},
\quad \text{soit }0,\quad \text{d'où}
\]
\[
(1)\qquad \chi(Y)=1-(\unsure{\gamma}+\mu)
\]
\margin{$Y$ surface compacte connexe sans bord non orientable de genre $\gamma$}

D'autre part, on \unsure{aura|a}
\[
\chi(X)=\chi(X\deleted{\unsure{)}}\setminus X^\sigma)
\]
car $X^\sigma$ est une \unsure{réunion} de cercles donc son $\chi$ est nul,

\scanpage{439}
\mspage{222}

Or, comme $X\setminus X^\sigma$ est un revêt. étale d'ordre $2$ de $\unsure{\operatorname{Int}Y}$, on a
\[
\chi(X\setminus X^\sigma)=2\chi(\operatorname{Int}(Y))
\]
enfin
\[
\chi(Y)=\chi(\underline{Y\setminus\partial Y})=\chi(\operatorname{Int}Y)
\]
pour la $\overline{\mathrm m}$ raison que \unsure{tantôt} $(\chi(\partial Y)=0)$,
d'où enfin
\[
(2)\qquad \chi(X)=2\chi(Y)=2(1-(\gamma+\mu))
\]
\deleted{\illegible{1}} ou encore, puisque $\chi(X)=2-2g$
\[
(3)\qquad g=\gamma+\mu
\]

Donc on trouve

\margin{\illegible{1}}

\uline{Proposition} La catégorie des surfaces orientées \inserted{connexes compactes} de genre $g$ munies d'une antiinvolution $\sigma$ est équivalente à celle des variétés compactes à bord \uline{non orientables}, de type $(\gamma,\mu)$ ($\gamma$ le genre, $\mu$ le nb des trous), avec $\gamma+\mu=g$.

\margin{non, on ne trouve qu'une sous-catégorie pleine de celle de tous les $(X,\sigma)$. Il y a d'autre part aussi les doublements orientés $Y\amalg_{\partial Y}Y$ des surfaces \uline{orientables} compactes à \uline{bord non vide} -- si $Y$ est de type $(\gamma,\mu)$, $X$ est de genre $g$ \unsure{avec} $g=2\gamma+\mu-1$ (où $\gamma\geq0$, $\mu\geq1$)}

\uline{Corollaire} Il y a exactement $g+1$ classes d'isomorphie de systèmes $(X,\sigma)$, classifiés par $0\leq\mu\leq g$, où $\mu=\operatorname{card}\unsure{\pi_0|\Pi_0}(X^\sigma)$.

\margin{Ça ne marche qu'en se limitant aux $(X,\sigma)$ tels que $(X/\sigma)$ non orientable -- \unsure{cf|cf.} ci-contre \deleted{\unsure{pour}} le cas \deleted{\unsure{général}} \inserted{$X/\sigma$ orientable}}

$X$ \inserted{$=X_g$} est une surface orientée compacte de

\deleted{Antiinvolutions des surfaces \unsure{orientées} de type \deleted{\illegible{1}} \unsure{algébroïde} (i.e. de la forme $X\setminus S$, avec $X$ compacte, $S$ fini). Nous \deleted{admettrons} \inserted{savons} qu'elles sont classifiées, mod isomorphisme \inserted{près}, par les systèmes}

\scanpage{440}

genre $g$, cela signifie aussi que dans le groupe $A_g=\mathrm{Aut}(X_g)$, il y a exactement $g+1$ classes de conjugaison d'éléments $\sigma$ satisfaisant
\[
(4)\qquad \sigma^2=1,\qquad \mathrm{sg}(\sigma)=+1.
\]
(i.e. qui soient des antiinvolutions). Elles fournissent $g+1$ classes de conjugaison d'éléments de $A_g/A_g^\circ=\unsure{\mathcal{T}_g}$, satisfaisant les \unsure{$\mathrm m^{\mathrm s}$|$\overline{\mathrm m}^{\mathrm s}$} relations. Nous montrerons (on l'espère!) que ces classes sont distinctes, et que tt $\sigma\in\unsure{\mathcal{T}_g}$ satisfaisant les relations (4) est dans l'une de ces classes.

On admettra le

\uline{Lemme} Toute antiinvolution du disque \inserted{unité} \inserted{(\unsure{cas|ou} de $\mathbb C$)} est conjuguée de $z\mapsto\overline z$, et par suite \unsure{l'ens.} de ses points fixes est homéomorphe à $\mathbb R$, et en particulier est \uline{connexe}.

\uline{Théorème} Soit $U$ une surface orientée \inserted{\unsure{(séparée connexe paracompacte)}}, \inserted{\uline{non isomorphe à} $S^2$,} $\sigma$ une antiinvolution de $U$, $\widetilde U$ un rev. universel de $U$, d'où $\pi=\mathrm{Aut}(\widetilde U/U)$ et une extension \unsure{$E$|$\mathcal E$} de \unsure{$\mathbb Z/2$} par $\pi$, formée
